\clearpage\SourceLeaf{196}%
\pdfbookmark[1]{Gauss an Schumacher, 1. April 1843}{schumacher18430401}
\fancyhead[L]{\small GAUSS AN SCHUMACHER}
\begin{center}
{\small GAUSS AN SCHUMACHER.\par}
\medskip\hrule\medskip
{\small Astronomische Nachrichten herausgegeben von H.\ C.\ \textsc{Schumacher}.\par
Band 20.\ Nr.\ 474.\ S.\ 299.\par}
\medskip\hrule
\end{center}\medskip
\begin{flushright}Göttingen, 1. April 1843.\end{flushright}
Um aus den Bahnelementen für eine gegebene Zeit einen Ort zu bestimmen, verwende ich zur Berechnung der Anomalie gern die Tafel von \textsc{Burckhardt}. Sie reicht aber nur bis $163^0\ 45'$ und genügt daher nach den Elementen von Herrn \textsc{Galle} für den gegenwärtigen Stand des Kometen nicht mehr. Die Tafel von \textsc{Barker} reicht zwar für alle Werte aus, ist aber bei großen Anomalien wegen der mühsamen Interpolation sehr unhandlich. In solchen Fällen verwende ich gewöhnlich ein besonderes Verfahren, das ich Ihnen mitteilen möchte und das Sie vielleicht interessiert. Ist $M$ die Zahl, mit der man (oder bei größeren Werten mit deren Logarithmen) in die Tafel von \textsc{Barker} eingehen müsste, also $M=\frac{\text{Zwischenzeit}}{nq^{\frac32}}$, wobei $\log n=0{,}0398723$, so setze ich $\log\frac{MM}{16875}=3P$ und versuche mithilfe meiner kleinen Logarithmentafel, wobei $A$ und $B$ die dort verwendete Bedeutung haben, die Gleichung $3A+2B=3P$ zu erfüllen. Das gelingt immer sehr schnell, wenn $P$ groß ist. Ist dann $a$ die zum Logarithmus $A$ gehörige Zahl und setzt man die Anomalie $=v$, so ergibt sich
\[
\operatorname{tang}\tfrac12v=\BareRoot{(3a)}\quad\text{oder}\quad
\log\operatorname{tang}\tfrac12v=\tfrac12(A+\log3)
\]

Auch der Logarithmus des Radiusvektors lässt sich dann. äußerst bequem berechnen. Man geht mit $A+\log3$ wieder in die erste Spalte ein, oder $A+\log3=A^*$, und die entsprechende Größe in der zweiten Spalte $=B^*$ folgt, womit sich sogleich der Logarithmus des Radiusvektors $=A^*+B^*+\log q$ ergibt.

\SourceLeaf{197}%
Zumindest für die ersten Versuche lässt sich jene Gleichung in der Form $C=P+\frac13B$ etwas bequemer und fast auf einen Blick indirekt lösen. Man kann zuerst $P$ in der dritten Spalte aufsuchen, oder $P=C'$ und die entsprechende Größe in der zweiten Spalte $=B'$ setzen, dann $P+\frac13B'=C''$ und dazu aus der Tafel die Größe der zweiten Spalte $=B''$ nehmen, dann (wo nötig) $P+\frac13B''=C'''$ und das dazugehörige $B'''$ nehmen und so weiter. Die Rechnung kommt sehr schnell zum Stillstand. Will man sich mit der Genauigkeit fünfstelliger Logarithmen nicht begnügen, so kann man hier mit Vorteil die Tafel von \textsc{Matthiessen} heranziehen (die ich sonst wegen der unangebrachten Sparsamkeit, mit der sie ganz unnötigerweise gedruckt ist, nicht gern verwende). Ich ziehe es aber vor, sie erst zu verwenden, wenn ich mithilfe der kleineren Tafel die beiden Stellen bestimmt habe, zwischen denen der endgültige Wert von $A$ liegt. Dann verwende ich lieber die Gleichung in ihrer ursprünglichen Form $3A+2B=3P$.

\Needspace{9\baselineskip}
Soll zum Beispiel die Anomalie für Februar $48.33333$ oder für die Zeit seit dem Periheldurchgang $20^{\mathrm T}87663$ bestimmt werden, so ergibt sich nach den Elementen von \textsc{Galle} [Astr.\ Nachrichten Band 20.\ N.\ 474 Berlin~1843~März~25]
\begin{center}
{\small\setlength{\tabcolsep}{3pt}
\begin{tabular}{r@{\ }r@{\qquad}r@{\ }r}
$q^{\frac32}\ .\ .\ .\ .\ .\ .$&$7.0809490$&$20.87663\ .\ .\ .\ .$&$1.3196604$\\
$n\BareRoot{16875}\ .\ .$&$2.1534942$&&$9.2344432$\\\cline{1-2}\cline{4-4}
Konst.\ Logarithmus $=$&$9.2344432$&&$2.0852172$\\[3pt]
&&Also $3P=$&$4.1704344$\\
&&$P=$&$1.3901448$
\end{tabular}}
\end{center}
Die kleinen Tafeln ergeben daraus
\[
\begin{array}{ll}
B'=0.01806&\quad C''=1.39616\\
B''=0.01781&\quad C'''=1.39608
\end{array}
\]
womit die Rechnung bereits feststeht und sich $A=1.37827$ ergibt. Die Tafel von \textsc{Matthiessen} gibt genauer $A=1.3782739$. Die weitere Rechnung verläuft dann so
\begin{center}
\begin{tabular}{r@{\ }l}
$A=$&$1.3782739$\\
$3\ .\ .\ .\ .\ .$&$0.4771213$\\\cline{2-2}
&$1.8553952$\\[3pt]
&$0.9276976=\log\operatorname{tang}\ 83^0\ 15'\ 49''53$\\
\multicolumn{2}{r}{und die wahre Anomalie $=\qquad166\ 31\ 39{,}06$}
\end{tabular}
\end{center}

\SourceLeaf{198}%
Ferner gehört zu
\begin{center}
\begin{tabular}{r@{\ }r}
$A^*=$&$1.8553952$\\
$B^*=$&$0.0060170$\\
$q\ .\ .\ .\ .$&$8.0539660$\\\cline{1-2}
Logarithmus des Radiusvektors $=$&$9.9153782$
\end{tabular}
\end{center}

Man sieht außerdem, dass diese Methode lediglich eine indirekte Lösung der bekannten kubischen Gleichung zwischen der Tangente der halben Anomalie und der Sektorfläche ist. Zugleich sieht man, dass sich meine Logarithmentafel oder, für eine genauere Rechnung, die von \textsc{Matthiessen} auf ganz ähnliche Weise dazu verwenden lässt, alle reellen Wurzeln \textit{jeder algebraischen Gleichung mit höchstens drei effektiven Gliedern} sehr bequem zu finden, wie ich kürzlich für die quadratische Gleichung in der jüngsten Ausgabe der Logarithmentafel von \textsc{Vega} gezeigt habe.
\par\medskip\centerline{\rule{.25\textwidth}{.3pt}}
